\documentclass[11pt,a4paper]{article}
\usepackage[T1]{fontenc}
\usepackage[utf8]{inputenc}
\usepackage{lmodern,amsmath,amssymb,amsthm,mathtools}
\usepackage[margin=25mm,headheight=16pt]{geometry}
\usepackage{microtype,booktabs,longtable,tabularx,array,enumitem}
\usepackage{xcolor,graphicx,fancyhdr,xurl}
\usepackage[colorlinks=true,linkcolor=navy,urlcolor=teal,citecolor=teal]{hyperref}
\definecolor{navy}{HTML}{16324F}
\definecolor{teal}{HTML}{087E8B}
\definecolor{pale}{HTML}{EEF4F7}
\hypersetup{pdftitle={Beyond Dirichlet GRH: Ideal Densities and Arithmetic Transfer Conditions},
 pdfauthor={Research report prepared for Vladimir Reshetnikov}}
\pagestyle{fancy}\fancyhf{}
\fancyhead[L]{\small\color{navy}Beyond Dirichlet GRH}
\fancyhead[R]{\small 26 November 2026}
\fancyfoot[C]{\thepage}
\setlength{\parskip}{5pt}\setlength{\parindent}{0pt}
\setlength{\emergencystretch}{2em}
\setlist{itemsep=3pt,topsep=4pt}
\newcolumntype{Y}{>{\raggedright\arraybackslash}X}
\newtheorem{theorem}{Theorem}[section]
\newtheorem{lemma}[theorem]{Lemma}
\newtheorem{proposition}[theorem]{Proposition}
\newtheorem{corollary}[theorem]{Corollary}
\theoremstyle{definition}\newtheorem{definition}[theorem]{Definition}
\theoremstyle{remark}\newtheorem{remark}[theorem]{Remark}
\newcommand{\Q}{\mathbb Q}\newcommand{\Z}{\mathbb Z}\newcommand{\N}{\mathbb N}
\newcommand{\F}{\mathbb F}\newcommand{\C}{\mathbb C}
\newcommand{\code}[1]{\texttt{\small #1}}
\newcommand{\upstream}[2]{\href{https://github.com/lsernau/grh-lean/blob/72fead102f663d5c7a407964a13dc0cf9d977207/#1}{#2}}
\newcommand{\incoming}[2]{\href{https://github.com/VladimirReshetnikov/ProveIt/blob/5cc267af75a635199ac5d93a7d49cbce521daefc/docs/incoming/#1}{#2}}

\begin{document}
\hypersetup{pageanchor=false}
\begin{titlepage}
\vspace*{12mm}
{\Large\color{teal}RESEARCH REPORT\par}
\vspace{9mm}
{\Huge\bfseries\color{navy}Beyond Dirichlet GRH\par}
\vspace{6mm}
{\LARGE Ideal densities and arithmetic\par transfer conditions\par}
\vspace{12mm}
{\large Prepared for Vladimir Reshetnikov\par}
{\large 26 November 2026\par}
\vspace{12mm}
\fcolorbox{navy}{pale}{\parbox{0.91\textwidth}{
\textbf{Main outcome.} We give an affirmative answer to the existence
part of Best et al.'s published Question 5.1, with an extension to
every number field: a finite change in prime-ideal valuation rules
strictly improves the natural density of the greedy set while
excluding all three-term geometric progressions with integral-ideal
ratios. We also prove strict separation between the optimal natural
and upper densities for both integral and fractional ideal ratios.

For the deeper arithmetic transfer, we construct compatible tame
characters over $\mathbb Q(\sqrt{-23})$, prove and check exact conductor
minima in a specified class, and correct a finite-precision objection
to the readable Lean source. These steps do not establish a new
nonabelian Riemann hypothesis or a new algebraic $K$-group vanishing.
}}
\vfill
{\small
\textbf{Frozen sources}\par
\texttt{lsernau/grh-lean}: \texttt{72fead102f663d5c7a407964a13dc0cf9d977207}\par
\texttt{VladimirReshetnikov/ProveIt}: \texttt{5cc267af75a635199ac5d93a7d49cbce521daefc}\par
Forty-six incoming GRH archives inventoried and their downloaded bytes checked.
\par\medskip
Human-readable proofs and exact finite certificates. New combined statements
are not presented as Lean-checked. Historical priority is not established.
}
\end{titlepage}
\hypersetup{pageanchor=true}

\begin{abstract}
The reusable content of the condensed RH/Dirichlet-GRH development is
more specific than its endpoint: a comparison of arithmetic covolumes,
prime-sensitive divisibility, almost-full spanning, controlled
anticyclotomic twists, and a one-sided zero detector. We review the
current incoming report corpus and pursue two different kinds of
transfer. First, extending its valuation-exchange method, we construct
progression-free integral ideals of natural density larger than the
greedy density in every number field. For $\Q(\sqrt{-23})$ the gain is
greater than $0.0011957$. A uniform prime-power obstruction then proves
that the supremum of upper densities exceeds the supremum of natural
densities for either integral-ideal or arbitrary fractional-ideal
ratios; both extrema are computable. Second, we construct auxiliary
characters satisfying the stronger precursor condition actually used
by the source, with arbitrarily small normalized tame cost, and certify
conductor minima at five budgets. Inspection of the readable source
also shows that its Selmer conclusion already has divisible
coefficients. We isolate a fixed-coefficient, single-precision
equivalent of its strict odd-cocycle input and identify the uniform
arithmetic estimates still needed for nonabelian transfer. The
combinatorial results use no RH assumption; the deeper transfer
remains incomplete.
\end{abstract}
\tableofcontents
\newpage

\section{What transfers, and what this report adds}

\subsection{The scope of the source theorem}

The five exported statements of the frozen repository concern the
Riemann zeta function and Dirichlet $L$-functions of all positive
moduli, including imprimitive characters.  Its half-plane statement
excludes zeros with real part greater than $1/2$, except for the
principal-character pole at $s=1$.  Neither the endpoint statement
nor its type quantifies over general Artin representations, general
Hecke characters of number fields, or all automorphic
$L$-functions~\cite{repo}.

We accept the supplied RH/Dirichlet-GRH result as the premise for
discussing its extensions. The source README reports standard Lean
axioms only and independent checker acceptance. This session did not
rebuild the full project: Lean and Lake were not installed. Its
theorem types, definitions, selected proof passages, and published
build claims were inspected separately from the computations replayed
here. The verification claims in this report refer to those
specific inspections and computations.

\subsection{The mechanism is a matched arithmetic comparison}

The source constructs a finite-dimensional space of target forms,
integral lattices, and chains of theta lifts. Write $D$ for its
dimension, $n$ for the number of independent short forms, $V$ for a
logarithmic squared covolume, and $\delta$ for an index cost.
In the source's own summary, the comparison has the form
\[
 H-\varepsilon_-\ \le\ \frac{V+2\delta}{D}
 \ \le\ H-\frac nD mR+\varepsilon_+,
\]
where $H$ is the same main height on both sides and $R$ is an
arithmetic prime-sum remainder. The error terms include the loss from
the missing rank; they cannot simply be discarded when replacing the
construction. The shared height cancels. In the original data the
resulting exact statement is
\[
                 m\,\mathrm{totalRemainder}
                      \le f a(\log b)^5.
\]
See \upstream{grh/src/New/Reduction/MainInequality.lean}
{\texttt{MainInequality.lean}} and the reading guide~\cite{guide}.

The three main inputs are short integral forms with controlled index,
enough independent forms, and a lower bound showing integral forms
are long. The prime-sensitive derivative estimates supply the
divisibility gain. Twists control congruence denominators in the
determinant lower bound. Finally a one-sided analytic criterion turns
the resulting prime estimate into zero exclusion. Reusing just the
last analytic criterion does not recreate the arithmetic estimates.

There are two useful meanings of ``transfer.'' A \emph{deep arithmetic
transfer} replaces the field, representation, and prime detector while
preserving these estimates. A \emph{component transfer} reuses an
ingredient, such as restrictions on exponents or specialization of a
power series, in a simpler problem. The complete ideal-density
results below are of the second kind. Their elementary antecedents
include Rankin and McNew~\cite{mcnew}; we do not attribute those classical ideas
to the RH proof.

\subsection{The current prior-work frontier}

The current ProveIt snapshot contains 46 relevant ZIP archives.
Comparison with the earlier snapshot \texttt{48efe71} finds five
additional archives and no modifications of the shared files. The
full inventory, exact Git blob hashes, SHA-256 checksums, and review
notes are included as data files. The review was targeted: all
archives were inventoried, while selected sections and checks were
read or replayed in depth. It was not a proof audit of every theorem
in every archive.

\begin{center}\small
\begin{tabularx}{\textwidth}{@{}>{\raggedright\arraybackslash}p{0.22\textwidth}YY@{}}
\toprule
Direction & What was already in the corpus & Increment in this report\\
\midrule
Progression-free sets & Polynomial and integer density constructions,
natural/upper gaps, finite optimization & Natural improvement and
strict natural/upper separation for ideals of every number field\\
Auxiliary twists & Ray-class obstruction and an order-three
conductor-$151^2$ example & Stronger precursor condition, all-depth
existence, five exact conductor minima\\
Selmer groups & Ramification-normalization and local-condition warnings
& Correction: full divisible bound already in readable source;
explicit descent for strict odd cocycles\\
Nonabelian zeros & Many equivalent positive prime-error criteria,
especially the discriminant-$-23$ $S_3$ example & Two arithmetic
interfaces clarified; the main prime-sensitive comparison remains open here\\
Algebraic $K$-theory & Bernoulli screens and proposed Hecke certificates
for degrees $124,128$ & Exact arithmetic prefilter replayed; no completed
large determinant or global cohomological bridge\\
Other applications & Class-group censuses, singular-modulus restrictions,
heavenly elliptic-curve claims, graph and finite-ring constructions
& Inventory and status distinctions; no blanket endorsement of
unreplayed claims\\
\bottomrule
\end{tabularx}
\end{center}

The previous reports already contain substantive results and claims.
It would be inaccurate to summarize the entire corpus as consisting
only of incomplete nonabelian reductions. Conversely, repeated
appearance of a result in several archives is not independent
verification. Sections~\ref{sec:ideal-natural-gain} and
\ref{sec:ideal-separation} give complete arguments for the extensions
established here.

\subsection{Answer to the original research question}

There is a concrete affirmative answer at the component level: the
valuation-exchange technique gives a theorem answering the existence
part of the published ideal-density Question 5.1 of Best
et al.~\cite{best}. It supplies new all-fields conclusions relative
to the inspected corpus, with an exact numerical certificate in the
same quadratic field used for the nonabelian benchmark.

For the deep arithmetic method, the nonabelian Hecke/Artin direction
remains a credible research program, with precise missing inputs.
We complete an auxiliary character construction and simplify the
odd-cocycle interface. We do not claim that these steps prove a new
nonabelian RH instance, Vandiver's conjecture, BSD, or the vanishing of
$K_{124}(\mathbb Z)$. A source-backed research program is useful only
when the unproved arithmetic estimates remain explicit.

\section{A natural-density improvement over the greedy ideals in every number field}
\label{sec:ideal-natural-gain}

Best, Huan, McNew, Miller, Powell, Tor and Weinstein ask whether the
greedy progression-free integral ideals in a quadratic number field can
be surpassed in natural density
(\href{https://arxiv.org/pdf/1412.0999}{Question~5.1 of their paper}).
The answer is affirmative in every number field. The argument extends
the finite valuation-cube exchange in the earlier
\href{https://github.com/VladimirReshetnikov/ProveIt/blob/5cc267af75a635199ac5d93a7d49cbce521daefc/docs/incoming/grh-progression-densities-and-cubic-prime-race.zip}{ProveIt polynomial-progression report}.
Its integer precedent is McNew's exchange construction. The new step
is to prove that a profitable finite exchange always exists over an
arbitrary number field, independently of its splitting patterns. The
proof uses no RH or GRH assumption. Historical priority is not asserted.

Let $K$ be a number field of degree $d$, let $\mathcal I_K$ be its nonzero
integral ideals, and put
\[
 \begin{aligned}
 A&=\left\{\sum_{j\ge0}\epsilon_j3^j:\epsilon_j\in\{0,1\},
                         \ \epsilon_j=0\text{ eventually}\right\},\\
 \mathcal G_K&=\{I\in\mathcal I_K:v_{\mathfrak p}(I)\in A
                         \text{ for every prime ideal }\mathfrak p\}.
 \end{aligned}
\]
Natural density is relative to the ideals of norm at most $X$.
A forbidden progression in this section is
$I,IJ,IJ^2$ with $J$ a proper nonzero \emph{integral} ideal.

\begin{theorem}\label{thm:ideal-natural-gain}
For every number field $K$, an explicitly decidable set
$\mathcal E\subseteq\mathcal I_K$ avoids all such progressions, has
natural density, and satisfies
\[
                         d(\mathcal E)>d(\mathcal G_K).
\]
There is an effective search for a finite set of prime ideals specifying
$\mathcal E$, with a positive rational certificate for the density gain.
\end{theorem}

\paragraph{The finite exchange.}
For a nonempty finite set $S$ of prime ideals, write $k=|S|$ and replace
the permitted local vectors $A^k$ by
\[
 D_S=\{u\in A^k:\min_i u_i=0\}
       \ \cup\ \bigl(\{1,2\}^k\setminus\{(1,\ldots,1)\}\bigr).
\]
Let $\mathcal E_S$ require this condition at $S$ and retain the condition
$v_{\mathfrak p}(I)\in A$ at every prime ideal outside $S$.

The set $A$ contains no nonconstant three-term arithmetic progression.
Indeed, in $a+c=2b$ the ternary digit sums have no carries, and a digit of
$2b$ is either zero or two. Equality forces the corresponding digits
of $a,b,c$ to agree. Therefore a progression of ideals in $\mathcal E_S$
would have zero valuation step outside $S$. Inside $S$ it would give
$u,u+h,u+2h\in D_S$ with $h\ge0$ and $h\ne0$.
If the final vector has a zero coordinate, all three vectors have a zero
there and lie in $A^k$, a contradiction. Otherwise the final vector is
in the new cube. At every moving coordinate its three entries must be
$0,1,2$. The initial vector consequently lies in $A^k$, so every fixed
coordinate belongs to $A\cap\{1,2\}=\{1\}$. The middle vector is the
deleted all-ones vector. This proves progression avoidance.

\paragraph{The density calculation.}
Set
\[
 w_{\mathfrak p}=(\operatorname N\mathfrak p)^{-1},\qquad
 T(w)=\sum_{\substack{a\in A\\a\ge1}}w^{a-1},\qquad
 p_{\mathfrak p}=(1-w_{\mathfrak p})\sum_{a\in A}w_{\mathfrak p}^{a},
 \qquad C_S=\prod_{\mathfrak p\notin S}p_{\mathfrak p}.
\]
All these products converge, and $C_S>0$. We use the classical ideal
counting asymptotic
$\#\{I:\operatorname NI\le X\}\sim\kappa_KX$, $\kappa_K>0$;
see \href{https://doi.org/10.1016/j.exmath.2006.08.001}{Murty--Van Order,
ideal-counting theorem}. Divisibility by an ideal $J$ has limiting proportion
$(\operatorname NJ)^{-1}$. Finite inclusion--exclusion gives independent
geometric valuation laws at any fixed finite set of prime ideals.
The unbounded conditions are justified by truncation: failure of the
greedy rule at a sufficiently large prime ideal requires
$\mathfrak p^2\mid I$. The counting asymptotic also gives a uniform
bound $\#\{I:\operatorname NI\le X\}\ll_K X$, so the upper density of
the omitted square-divisibility events is
$O_K(\sum_{\operatorname N\mathfrak p>Y}(\operatorname N\mathfrak p)^{-2})$,
which tends to zero. Large valuations at finitely many remaining primes
have a uniformly vanishing geometric tail. Thus both densities exist,
and subtraction of the deleted positive orthant and added cube gives
\begin{equation}\label{eq:ideal-natural-delta}
 d(\mathcal E_S)-d(\mathcal G_K)
 =C_S\prod_{\mathfrak p\in S}(1-w_{\mathfrak p})w_{\mathfrak p}
 \left[\prod_{\mathfrak p\in S}(1+w_{\mathfrak p})-1
                    -\prod_{\mathfrak p\in S}T(w_{\mathfrak p})\right].
\end{equation}

\paragraph{Why some finite exchange is profitable.}
The support of $T$ is contained in $\{0\}\cup\{2,3,\ldots\}$.
Every integer at least two is $2a+3b$ for some $a,b\ge0$, whence
\[
 T(w)\le\frac1{(1-w^2)(1-w^3)},\qquad
 \prod_{\mathfrak p\in S}T(w_{\mathfrak p})
       \le\zeta_K(2)\zeta_K(3)
       \le\bigl(\zeta(2)\zeta(3)\bigr)^d<2^d.
\]
The middle inequality follows prime by prime: at most $d$ prime ideals
lie over a rational prime, and each has norm at least that prime.
For completely rational bounds, integral tail estimates give
\[
 \zeta(2)<\sum_{n=1}^{10}n^{-2}+\frac1{10}
    =\frac{2095337}{1270080}<\frac{33}{20},\qquad
 \zeta(3)<\sum_{n=1}^{4}n^{-3}+\frac1{32}
    =\frac{2089}{1728}<\frac{121}{100};
\]
their last two bounds multiply to $3993/2000<2$.
On the other hand, $\sum_{\mathfrak p}(\operatorname N\mathfrak p)^{-1}$
diverges: otherwise the Euler product at $1$ would converge, contrary
to the ideal-counting asymptotic. Thus
$\prod_{\mathfrak p\in S}(1+w_{\mathfrak p})$ is unbounded along
increasing finite sets $S$.

Consequently the finite rational test
\begin{equation}\label{eq:ideal-natural-certificate}
                  \prod_{\mathfrak p\in S}(1+w_{\mathfrak p})>1+2^d
\end{equation}
eventually succeeds. Moreover
$C_S\ge\zeta_K(2)^{-1}>2^{-d}$, because $0,1\in A$.
The expression obtained from the right side of
\eqref{eq:ideal-natural-delta} by replacing $C_S$ by $2^{-d}$ and
$\prod T(w_{\mathfrak p})$ by $2^d$ is therefore a positive rational
lower bound for the gain whenever \eqref{eq:ideal-natural-certificate}
holds. Enumerating and factoring rational primes gives a terminating
search for $S$; factoring an input ideal and testing its valuation
digits decides membership in $\mathcal E_S$. This proves the theorem.

\paragraph{A short certificate over $\mathbb Q(\sqrt{-23})$.}
Write $\omega=(1+\sqrt{-23})/2$, whose minimal polynomial is
$X^2-X+6$. At each of $p=2,3$, its two distinct roots modulo $p$ are
$0,1$. The four prime ideals $(p,\omega)$ and $(p,\omega-1)$ therefore
have norms $2,2,3,3$. Take these four ideals for $S$.
The sharper bound
\[
                T(w)\le\overline T(w):=1+w^2+w^3+\frac{w^8}{1-w}
\]
follows from $A\cap[1,8]=\{1,3,4\}$. Exact arithmetic gives
\[
 \overline T(1/2)=\frac{177}{128},\qquad
 \overline T(1/3)=\frac{5023}{4374},\qquad
 \prod_{\mathfrak p\in S}(1+w_{\mathfrak p})-1
       -\prod_{\mathfrak p\in S}\overline T(w_{\mathfrak p})
       =\frac{16658080679}{34828517376}.
\]
Also $\prod_{\mathfrak p\in S}(1-w_{\mathfrak p})w_{\mathfrak p}=1/324$.
Since $\prod_{\mathfrak p\in S}(1-w_{\mathfrak p}^2)=4/9$ and
$\zeta_K(2)\le\zeta(2)^2<(5/3)^2$, we have
\[
 C_S\ge\prod_{\mathfrak p\notin S}(1-w_{\mathfrak p}^2)
       =\frac{\zeta_K(2)^{-1}}{4/9}>\frac{81}{100}.
\]
The resulting explicit natural-density gain is
\begin{equation}\label{eq:ideal-natural-23-gain}
 d(\mathcal E_S)-d(\mathcal G_K)
 >\frac{16658080679}{13931406950400}>0.0011957.
\end{equation}
The supplied rational checker verifies the prime reductions and every
number in this certificate. It does not substitute for the all-fields
proof above.

\paragraph{The precise boundary.}
The theorem treats integral-ideal ratios. The displayed exchange permits
mixed-step triples: for two selected coordinates, $(0,3),(1,2),(2,1)$
all satisfy the new rule, and further coordinates may be fixed at one.
Accordingly this construction does not answer Best et al.'s Question~5.4
about a \emph{natural}-density improvement for fractional-ideal ratios.
It also does not determine the optimal natural density. Those conclusions
must be separated from a strict natural-versus-upper-density theorem.

\section{Natural and upper densities separate in every number field}
\label{sec:ideal-separation}

This section extends the smooth-factor optimization in the incoming
\path{grh_transfers_density_gaps_and_zero_tests.zip}~\cite{prior-density-gaps} from integers
and polynomials to ideals of an arbitrary number field.  Its additional
ingredient is a single prime-power obstruction that works in every field,
including when distinct ideals have the same norm.  All assertions in this
section are unconditional.

Let $K$ be a number field of degree $d$, and let $\mathcal I_K$ be its
nonzero integral ideals.  Put
\[
 A_K(X)=\#\{\mathfrak a\in\mathcal I_K:N\mathfrak a\le X\}.
\]
We use the classical ideal-counting theorem
$A_K(X)=\kappa_K X+o(X)$, where $\kappa_K>0$; see
\cite{murty-vanorder}.  Natural density means a limit of the
proportion in this norm ball, and upper density means its limsup.

Fix either of two conventions.  In the \emph{integral} convention the
forbidden triples are
$\mathfrak a,\mathfrak a\mathfrak r,
\mathfrak a\mathfrak r^2$ with integral
$\mathfrak r\ne\mathcal O_K$.  In the \emph{fractional} convention
$\mathfrak r\ne\mathcal O_K$ may be any nonzero fractional ideal, as long
as all three terms are integral.  The latter includes nontrivial ratios
of norm one.  In valuation coordinates these are, respectively,
nonconstant arithmetic progressions with a nonnegative step and with an
arbitrary integral step.  Write $b_K$ for the supremum of natural
densities of admissible sets, and $u_K$ for the supremum of upper
densities, using the same fixed convention in both definitions.

\begin{theorem}[Separation and finite approximation]
\label{thm:ideal-separation}
For every number field and for either convention,
\[
                         u_K>b_K.
\]
Both numbers are computable from an explicit presentation of $K$.
More precisely, fix a prime ideal $\mathfrak p$, put
$Q=N\mathfrak p$ and $M=Q^8$.  For every integer $Y\ge Q$, define
\[
 c_Y=\prod_{N\mathfrak q\le Y}(1-(N\mathfrak q)^{-1}).
\]
Then
\begin{equation}\label{eq:ideal-gap}
 u_K-b_K\ \ge\ \frac{c_Y}{M(M+1)}-\frac{d}{Y}.
\end{equation}
The right side is positive for every sufficiently large $Y$.
\end{theorem}

\subsection{The two local optimization problems}

Let $S$ be the prime ideals of norm at most $Y$.  An ideal is
$S$-smooth if it has no prime factor outside $S$, and $S$-rough if it
has none in $S$.  Every ideal has a unique factorization
$\mathfrak a=\mathfrak g\mathfrak h$ with rough $\mathfrak g$ and
smooth $\mathfrak h$.  Set
\[
 c=\prod_{\mathfrak q\in S}(1-(N\mathfrak q)^{-1}),\qquad
 v=\prod_{\mathfrak q\notin S}(1-(N\mathfrak q)^{-2}).
\]
For admissible subsets $D$ of the smooth ideals, and for $t\ge1$, put
\begin{align}
 B&=c\sup_D\sum_{\mathfrak h\in D}\frac1{N\mathfrak h},\\
 J(t)&=\max\{\#D:D\text{ admissible},\ N\mathfrak h\le t
                         \text{ for all }\mathfrak h\in D\},\\
 U&=c\int_1^\infty J(t)\,\frac{dt}{t^2}.
\end{align}
The smooth reciprocal sum is exactly $1/c$.  In particular all these
expressions converge, $B\le U\le1$, and the weighted supremum defining
$B$ is attained.  For the latter assertion one may use compactness of
the space of subsets of the countable smooth monoid: admissibility is
closed and the summable weighted objective is continuous.

\begin{lemma}[Local optimization bounds]\label{lem:BU}
For either ratio convention,
\[
                  vB\le b_K\le B,\qquad vU\le u_K\le U.
\]
\end{lemma}
\begin{proof}
Finite inclusion--exclusion and ideal counting give rough ideals
relative density $c$.  Squarefree rough ideals have relative density
$cv$.  Indeed, impose squarefreeness at finitely many further primes
and then let that finite set increase.  The omitted-square error is
bounded by a constant times $\sum(N\mathfrak q)^{-2}$, uniformly as
$X\to\infty$, since $A_K(t)\ll_K t$.  This convergent bound justifies
passage to the infinite product.

For the upper bound on $b_K$, let $E$ have natural density $b$.  Partial
summation yields
\[
 \sum_{\substack{\mathfrak a\in E\\N\mathfrak a\le X}}
       \frac1{N\mathfrak a}
   =(b\kappa_K+o(1))\log X.
\]
Each rough fiber $D_{\mathfrak g}$ is admissible.  Consequently its
smooth reciprocal sum is at most $B/c$, and the full sum is bounded by
\[
 \frac Bc\sum_{\substack{\mathfrak g\text{ rough}\\N\mathfrak g\le X}}
          \frac1{N\mathfrak g}
   =(B\kappa_K+o(1))\log X.
\]
Thus $b\le B$.

For the upper bound on $u_K$, count in each rough fiber at cutoff $X$:
\[
 \#\{\mathfrak a\in E:N\mathfrak a\le X\}
 \le \sum_{\substack{\mathfrak g\text{ rough}\\N\mathfrak g\le X}}
                  J(X/N\mathfrak g).
\]
Rough ideal counting, first on a bounded interval $1\le t\le T$, gives
limit $c\int_1^T J(t)t^{-2}\,dt$ after division by $A_K(X)$.
The tail is uniformly controlled by the corresponding smooth counting
function $H(t)$: $H(t)=O_S((1+\log t)^{|S|})$ and
$\int_1^\infty H(t)t^{-2}\,dt=1/c$.  Equivalently, expand by smooth
ideals and dominate each term by $O_K(1/N\mathfrak h)$.
Letting $T\to\infty$ proves the upper bound.

For both lower bounds retain only squarefree rough parts.  In any
progression, each valuation outside $S$ would form an arithmetic
progression in $\{0,1\}$, so it is constant.  All three terms must
therefore lie in a single rough fiber.  Using the same admissible
smooth set $D$ in every such fiber gives natural density
$cv\sum_{\mathfrak h\in D}1/N\mathfrak h$.  Taking the supremum
proves $b_K\ge vB$.

To attain the lower upper-density bound, choose $X_j$ with
$X_j/X_{j-1}\to\infty$.  For every squarefree rough $\mathfrak g$
with $X_{j-1}<N\mathfrak g\le X_j$, permanently assign a maximizing
smooth set through $X_j/N\mathfrak g$.  This is one infinite admissible
set, since each fiber is assigned once.  At cutoff $X_j$ and fixed
$T$, its newly assigned fibers with $X_j/T<N\mathfrak g\le X_j$
give limiting proportion $cv\int_1^T J(t)t^{-2}\,dt$.
Let $T\to\infty$ to obtain upper density at least $vU$.
\end{proof}

\subsection{A uniform obstruction to compatible prefix optima}

Let $A$ be the nonnegative integers with ternary digits only $0,1$.
The greedy smooth set $G$ consists of ideals all of whose prime
exponents lie in $A$.  The following proof also resolves the issue of
equal norms: a norm layer need not be either entirely included or
entirely excluded.

For $n=\sum_j d_j3^j$, define
\[
 a(n)=\sum_{d_j=1}3^j,\qquad
 b(n)=\sum_{d_j\in\{1,2\}}3^j.
\]
Then $a(n),b(n)\in A$ and $n=2b(n)-a(n)$.  If $n\in A$ both equal
$n$; otherwise $0\le a(n)<b(n)<n$.  Applying this coordinatewise
shows that every nongreedy ideal is the final term of a progression
whose first two terms are greedy ideals of strictly smaller norm.
The ratio is integral.  Meanwhile all greedy ideals coexist without
any progression, even with fractional ratios.

It follows by induction on integer norm that any admissible set
simultaneously maximizing the number of ideals at \emph{every} cutoff
$1,2,\ldots,M$ must equal $G$ through $M$.  At a new norm layer every
nongreedy candidate is already blocked, and all greedy candidates can
be included together.  To achieve even the greedy count one must
include every such candidate.  This argument does not choose an
ordering within a norm layer.

Now change only the pure powers of the selected prime $\mathfrak p$
at cutoff $M=Q^8$.  Replace exponent set
\[
                    \{0,1,3,4\}\quad\text{by}\quad
                    \{0,1,3,7,8\}.
\]
The latter is arithmetic-progression-free and has one additional
element.  The resulting smooth set is admissible even for fractional
ratios.  To check this, any new progression has an added pure
$\mathfrak p$-power as one term.  At every other prime one of its
three valuations is zero; the other two belong to $A$ (or are also
zero).  Since $0\in A$ and $A$ is progression-free, all three are
zero.  The progression is confined to pure $\mathfrak p$-powers,
where the new exponent set forbids it.

Thus $G$ is not cardinality-optimal through $M$.  No admissible set
can optimize all prefixes through $M$.  For every admissible smooth
$D$ there is some integer $1\le k\le M$ with
$J(k)-\#\{\mathfrak h\in D:N\mathfrak h\le k\}\ge1$.
That deficit persists throughout $[k,k+1)$, giving
\[
 \int_1^\infty
 \bigl[J(t)-\#\{\mathfrak h\in D:N\mathfrak h\le t\}\bigr]
          \frac{dt}{t^2}
 \ge\frac1{k(k+1)}\ge\frac1{M(M+1)}.
\]
Taking the infimum over $D$ proves
\begin{equation}\label{eq:BUgap}
                         U-B\ge\frac{c}{M(M+1)}.
\end{equation}

\subsection{Passing to the global densities}

There are at most $d$ prime ideals of any specified integer norm:
that norm must be a power of a unique rational prime, and at most $d$
primes lie over it.  Therefore
\[
 1-v\le\sum_{N\mathfrak q>Y}(N\mathfrak q)^{-2}
       \le d\sum_{n>Y}n^{-2}\le d/Y.
\]
Equations~\eqref{eq:BUgap} and Lemma~\ref{lem:BU} now give
\[
 u_K-b_K\ge vU-B\ge U-B-(1-v)
                     \ge c_Y/[M(M+1)]-d/Y.
\]
Finally
\[
 c_Y\ge\prod_{p\le Y}(1-p^{-1})^d\gg_d(\log Y)^{-d}
\]
by the classical Mertens product theorem.  This decreases more slowly
than $d/Y$, proving strict positivity.  One could instead use the
number-field Mertens theorem, but it is unnecessary here.

For computability choose $Y$ with $d/Y$ below the desired tolerance.
The finite smooth reciprocal sum through $H$ approaches $1/c$, with
explicit rational tail
$1/c-\sum_{N\mathfrak h\le H}1/N\mathfrak h$.
Finite weighted independent-set optimization through $H$ approximates
$B$ with at most $c$ times this tail.  Optimizing cardinalities at
each norm through $H$ and integrating the resulting step function
approximates $U$; the tail is bounded by
$c\int_H^\infty H(t)t^{-2}\,dt$.  Writing $H_S(t)$ for the smooth
counting function to avoid confusion with the cutoff, this integral is
the computable rational number
\[
 c\left(\frac{H_S(H)}{H}+\frac1c
             -\sum_{N\mathfrak h\le H}\frac1{N\mathfrak h}\right).
\]
It tends to zero as $H$ increases. Lemma~\ref{lem:BU} then
computes $b_K$ and $u_K$.  This is an existence-of-algorithm result;
the exact independent-set problems can be expensive.

\subsection{An explicit certificate in the conductor-23 field}

For $K=\mathbb Q(\sqrt{-23})$, choose a prime of norm $2$, so $M=256$.
The script \path{code/check_ideal_density.py} computes $c_Y$ at
$Y=10^7$ using only integer arithmetic.  For each rational prime,
the splitting test is its quadratic-residue class modulo $23$;
$23$ ramifies, split primes contribute two factors, and inert primes
contribute one factor when $p^2\le Y$.

After each Euler factor the script rounds \emph{down} on the grid
$10^{-40}\mathbb Z$.  Its certified lower bound is
\[
 c_{10^7}\ge
 \frac{11078476589752147709803027083658225201}
      {625000000000000000000000000000000000000}.
\]
There are $664449$ factors.  Using the exact tail bound $2/10^7$
in~\eqref{eq:ideal-gap} proves, under either convention,
\begin{equation}\label{eq:ideal-explicit-gap}
                         u_K-b_K>\frac1{20000000}.
\end{equation}
The sharper rational lower bound is stored in the JSON certificate.
This computation checks a finite numerical part of the theorem; the
norm-counting limits and infinite construction are established by the
proof above.

For fractional ratios, the theorem implies upper density strictly
larger than greedy density, since the greedy set has a natural
density at most $b_K$.  It does \emph{not} show that the fractional-ratio
natural-density supremum exceeds greedy density.  Thus it contributes
to the extensions requested in Questions 5.3 and 5.5 of
Best et al.~\cite{best}, while their Question 5.4 remains unresolved
by this report.  No assertion of priority for the general separation
theorem is made.

\section{Realizing the actual auxiliary character condition}
\label{sec:source_twists}

The previous bridges report identified an obstruction to extending local
torsion values to global characters and computed the minimum conductor
$151^2$ for a cubic anticyclotomic character over
$K=\mathbb Q(\sqrt{-23})$ that is trivial above $3$.
\cite{prior-bridges}
There is a further distinction in the source.  Its predicate
\texttt{IsFromU1} asks for a character of the form
\[
                 \psi=\xi/\xi^c,
\]
where the precursor $\xi$ is ramified on only one side of each conjugate
pair and is itself trivial on the prescribed local components.
Triviality of $\psi$ alone does not imply this stronger assertion.
The source explicitly describes both the precursor and its triviality on
the whole decomposition groups, rather than only their inertia subgroups.
\cite{sourceCMSplit,sourceAuxPrimes}

This section makes that distinction effective in the discriminant-$-23$
benchmark.  It gives exact conductor minima for five small valuation
budgets and an existence construction for every positive budget.
These are character-theoretic results; the Selmer and determinant
requirements are additional conditions.

\subsection{A precise optimization problem}

Let $c$ be complex conjugation and let $S$ consist of the two primes of
$K$ above $3$.  Choose a finite set $Q$ of primes of prime norm, disjoint
from $S$, with $Q\cap cQ=\varnothing$.  Thus each rational prime in the
support splits in $K$, and $Q$ contains one of its two primes.
We consider nontrivial continuous finite-order characters
\[
  \xi:G_K\longrightarrow\mu_{3^\infty},\qquad
  \psi=\xi/\xi^c,
\]
where $\xi$ is tamely ramified only on $Q$ and is trivial on the
decomposition groups at both members of $S$.  Delete any primes of $Q$
where $\xi$ is unramified, so that $Q$ is the actual support.
Write
\[
  \mathcal T(\psi)=
  \sum_{\mathfrak v\mid\mathfrak f(\psi)}
       v_3\bigl(1-\psi(\tau_{\mathfrak v})\bigr),
       \qquad v_3(3)=1,
\]
with $\tau_{\mathfrak v}$ a tame inertia generator.  This quantity is
independent of those choices.  If the local order at pair $i$ is
$3^{n_i}$, then
\begin{equation}\label{eq:source_twist_cost}
 \mathcal T(\psi)=\sum_i3^{1-n_i},\qquad
 \operatorname N\mathfrak f(\psi)=\left(\prod_iq_i\right)^2.
\end{equation}
Indeed, each primitive $3^{n_i}$th root contributes
$1/(2\cdot3^{n_i-1})$ to the valuation, and the pair contributes twice
that number.  The valuation formula follows from
$\Phi_{3^{n_i}}(1)=3$ and equality of the valuations of its conjugate
factors.  Tameness makes each nontrivial conductor exponent one.

\begin{theorem}[Exact conductor minima with the precursor condition]
\label{thm:source_twist_minimum}
For $1\le N\le5$, impose the budget
$\mathcal T(\psi)\le3^{1-N}$.  Among all the characters just specified,
including characters supported on several split pairs, the exact minimum
of $\operatorname N\mathfrak f(\psi)$ and the number of output
characters attaining it are as follows.
\begin{center}
\begin{tabular}{rcrrr}
\hline
$N$ & Budget & $q_{\min}$ & Minimum conductor norm & Number\\
\hline
1 & $1$    & 151       & 22801             & 2\\
2 & $1/3$  & 6481      & 42003361          & 24\\
3 & $1/9$  & 6481      & 42003361          & 18\\
4 & $1/27$ & 11767033  & 138463065623089   & 54\\
5 & $1/81$ & 102906613 & 10589770999131769 & 162\\
\hline
\end{tabular}
\end{center}
Every minimum is attained on one conjugate pair, with conductor ideal
$(q_{\min})$ for $\psi$.  For budget $1/3$ the minimum includes six
characters with local order $9$ and eighteen with local order $27$.
\end{theorem}

\subsection{The exact ray-group reduction}

Put
\[
 \theta=(1+\sqrt{-23})/2,\qquad
 \beta=1+2\theta=2+\sqrt{-23},\qquad
 \mathfrak s=(3,\theta-1).
\]
The three reduced primitive forms of discriminant $-23$ are
$(1,1,6),(2,1,3),(2,-1,3)$; the reduced-form bound
$a\le\sqrt{23/3}<3$ makes the list exhaustive.  Thus
$\operatorname{Cl}(K)\simeq C_3$.
The norm form $u^2+uv+6v^2$ does not represent $3$, so
$\mathfrak s$ is nonprincipal and generates the class group.
Since $\operatorname N\beta=27$ and $\beta$ is a unit at the other
prime over $3$,
\begin{equation}\label{eq:source_twist_class_relations}
       \mathfrak s^3=(\beta),\qquad
       \mathfrak s\,c\mathfrak s=(3).
\end{equation}
The $S$-units are exactly $\{\pm3^a\beta^b:a,b\in\mathbb Z\}$:
their valuation vectors $(u,v)$ satisfy $u-v\equiv0\pmod3$, and
$(1,1),(3,0)$ generate precisely that lattice.

Let
\[
 R_Q=\prod_{\mathfrak q\in Q}
       (\mathcal O_K/\mathfrak q)^\times[3^\infty].
\]
The ordinary ray exact sequence has kernel $R_Q$ on the $3$-primary
part, because the units of $K$ are $\{\pm1\}$.  Its quotient is $C_3$.
A lift $b$ of $[\mathfrak s]$ satisfies $3b=r_\beta$, and the class
of $c\mathfrak s$ is $r_3-b$, where $r_u$ denotes the local residue
tuple of $u$, projected to $R_Q$.  A consistent inversion convention
for principal ideals changes signs but not the subgroups below.
Consequently, imposing triviality on \emph{both} decomposition groups
above $3$ gives the exact quotient
\begin{equation}\label{eq:source_twist_quotient}
 \frac{\operatorname{Cl}_Q(K)[3^\infty]}
 {\langle[\mathfrak s],[c\mathfrak s]\rangle}
       \simeq R_Q/\langle r_3,r_\beta\rangle.
\end{equation}
The characters of this quotient are exactly the allowed precursors
$\xi$.  This is also the ray class sequence with the $S$-units
quotiented out; the classical class-field correspondence identifies it
with the stated Galois-character conditions.\cite{milneCFT}

For a single split rational prime $q$, put $h=v_3(q-1)$ and choose a
root $r$ of $r^2=-23$ in $\mathbb F_q$.  The local group is cyclic
of order $3^h$.  A precursor with local order at least $3^N$ exists
if and only if
\begin{equation}\label{eq:source_twist_residue_tests}
 \begin{split}
  &q\equiv1\pmod{3^N},\qquad \left(\frac{-23}{q}\right)=1,\\
  &3^{(q-1)/3^N}\equiv1\pmod q,
       \qquad(2+r)^{(q-1)/3^N}\equiv1\pmod q.
 \end{split}
\end{equation}
These tests do not depend on which square root is chosen: the other
value of $\beta$ is $27/\beta$, and the test on $3$ already controls
$27$.

\subsection{Proof and exact verification of the minima}

The verifier \texttt{check\_source\_twists.py} lists every prime in
the progression $1\bmod3^N$ up to the candidate bound using an
Eratosthenes sieve of that progression, then applies
\eqref{eq:source_twist_residue_tests}.  It uses no external packages,
probable-prime tests, numerical approximations, or class-group oracle.
Its exact search finds the displayed first primes.  The exponents of
the cyclic quotients in \eqref{eq:source_twist_quotient} at those primes
are respectively $1,3,3,4,5$.

Here is why this finite search also proves the minimum over arbitrary
supports.  If there are at least two pairs and the budget is
$3^{1-N}$, each $n_i$ must be at least $N+1$: a single term with
$n_i\le N$ would already exhaust or exceed the budget, leaving no
room for another positive term.  Hence every support prime is split
and is $1\bmod3^{N+1}$.  The first two such primes are
\[
\begin{array}{c|rrrrr}
N&1&2&3&4&5\\ \hline
q_1&73&163&163&487&42283\\
q_2&127&271&487&3889&52489.
\end{array}
\]
For $N=1,2,3,5$, their product exceeds the corresponding
single-support minimum $q_{\min}$, which excludes every multipair
support.

For $N=4$, three pairs cost at least $487^3>11767033$ before the
conductor is squared.  Since $3889^2>11767033$, every potentially
cheaper two-pair support must be $(487,q)$, with
\[
                q\in\{3889,4861,8263,19441,19927\}.
\]
All these primes have $v_3(q-1)=5$.  Thus both local characters
would have to have exact order $3^5$.  Choose generators of their
local cyclic $3$-groups and reduce the coordinate vectors of
$(r_3,r_\beta)$ modulo $3$.  One exact choice gives
\[
\begin{array}{c|rrrrrr}
q&487&3889&4861&8263&19441&19927\\ \hline
\log(3)\bmod3&1&0&2&1&0&0\\
\log(\beta)\bmod3&1&2&0&0&1&2.
\end{array}
\]
The column for $487$ is independent of every other column, even
after replacing the second row in that other column by its negative.
The sign change covers the other choice of the prime above $q$,
since $\log(\bar\beta)=3\log(3)-\log(\beta)$.
Thus neither orientation admits nonzero local character coefficients
modulo $3$ satisfying both relations.  This excludes all five pairs.

Finally, at a cyclic quotient $C_{3^k}$, the number of characters
with local order at least $3^N$ is
$\sum_{j=N}^{k}\varphi(3^j)$, which gives the final column of the
table.  The map $\xi\mapsto\xi/\xi^c$ is injective for a fixed
half-support: its kernel would be unramified everywhere, while
triviality at $\mathfrak s$ kills every unramified character because
$[\mathfrak s]$ generates the class group.  Choosing the other half
does not add new outputs, since it replaces a precursor by the inverse
of its conjugate.  This completes the proof of
Theorem~\ref{thm:source_twist_minimum}.

\subsection{Why the stronger condition changes the answer}

For comparison, if only the anticyclotomic output $\psi$ is required
to be trivial above $3$, the one-pair quotient is
\[
 C_{3^h}/\langle r_{\beta/\bar\beta}\rangle.
\]
The corresponding exact minimal support primes at budgets
$1,1/3,1/9,1/27,1/81,1/243$ are
\[
                  151,\ 541,\ 6481,\ 94447,\ 238627,\ 1573183.
\]
The same verifier checks their minima, including smaller multipair
supports.  These are \emph{not} substitutes for the stronger table.
For example, at $q=541$, choose $r=40$, so $\beta=42$.  Then
\[
 \left(\frac{2+r}{2-r}\right)^{60}=1\pmod{541},\qquad
             3^{60}=352\pmod{541},\qquad
            42^{60}=129\pmod{541}.
\]
The weaker character exists at budget $1/3$, but the precursor
condition fails.  Enforcing the actual precursor condition raises
the minimum conductor norm from $292681$ to $42003361$.

\subsection{Arbitrarily small budgets and chains of powers}

\begin{theorem}[A governing field for the precursor condition]
\label{thm:source_twist_all_budgets}
For each $N\ge1$, set $k=3^N$ and
\[
 L_N=K\bigl(\zeta_k,3^{1/k},\beta^{1/k}\bigr).
\]
This is a finite Galois extension of $\mathbb Q$, of degree at most
$4\cdot3^{3N-1}$.  Every rational prime splitting completely in
$L_N$ supplies a precursor $\xi$ as above, with local order $3^N$
on one chosen prime above $q$.  Its output $\psi=\xi/\xi^c$ has
conductor $(q)$ and $\mathcal T(\psi)=3^{1-N}$.

More generally, given an integer $m\ge1$ and $\varepsilon>0$, choose
$N>\lfloor\log_3m\rfloor$ such that
\[
             3^{1-N+\lfloor\log_3m\rfloor}\le\varepsilon.
\]
Then a character exists for which the first $m$ positive powers all
have conductor $(q)$ and satisfy
$\mathcal T(\psi^j)\le\varepsilon$ for $1\le j\le m$.
The support can avoid any prescribed finite set of rational primes.
\end{theorem}

\begin{proof}
The field contains all conjugates of the displayed radicals.
In particular, $c\beta=27/\beta$, so a $k$th root of $c\beta$ is
$(3^{1/k})^3/\beta^{1/k}$, and the other roots differ by elements
of $\mu_k$.  Thus $L_N/\mathbb Q$ is Galois.  Adjoining the two
radicals to $K(\zeta_k)$ gives
\[
       [L_N:\mathbb Q]\le2\varphi(k)k^2
                           =4\cdot3^{3N-1}.
\]
Complete splitting makes $q$ split in $K$, gives $q=1\bmod k$,
and makes $3$ and $\beta$ $k$th powers in the residue field.
The tests \eqref{eq:source_twist_residue_tests} therefore hold.
The quotient in \eqref{eq:source_twist_quotient} has a character
with exact local order $k$, giving the claimed precursor.
Chebotarev's theorem gives infinitely many completely split primes,
so any finite set may be excluded.  This use is unconditional and
qualitative.\cite{milneCFT}

For $1\le j\le m<3^N$, the inertia value of $\psi^j$ has order
$3^{N-v_3(j)}$.  Hence
\[
        \mathcal T(\psi^j)=3^{1-N+v_3(j)}
            \le3^{1-N+\lfloor\log_3m\rfloor}.
\]
No such power loses ramification, proving the last assertions.
\end{proof}

For $0<\varepsilon\le1$, choosing $N$ minimally in this theorem
also gives the explicit degree estimate
\[
            [L_N:\mathbb Q]<972\,(m/\varepsilon)^3.
\]
This bounds the degree of the field governing the prime-selection
problem.  It does not bound the least completely split prime or the
conductor.  In particular, a least-prime estimate for rational
arithmetic progressions cannot be applied to $L_N$ merely from
Dirichlet GRH.

\subsection{What has and has not been transferred}

The upstream proof already uses roots of $S$-units and Chebotarev to
produce auxiliary characters.\cite{sourceAuxPrimes}
The contribution here is the explicit quotient
$R_Q/\langle3,\beta\rangle$, the exact conductor optimization,
the distinction between the output and precursor conditions, and
the degree estimate for this fixed quadratic benchmark.
The five-budget precursor optimization and its exact certificates
extend the reviewed incoming packages, which already contain the
order-three conductor-$151^2$ example. No literature-priority claim
is made.

These characters satisfy the local precursor, tame-ramification,
decomposition-group, and small-valuation requirements isolated here.
They do not establish the simultaneous Selmer length estimates,
special-value bounds, or prescribed Frobenius neighborhoods needed
for regular twists in the determinant argument.  Also,
$K=\mathbb Q(\sqrt{-23})$ is a benchmark base field, not the full
cyclotomic CM field of the upstream admissible data; a base extension
requires recomputing its units and local conditions.  No nonabelian
Riemann hypothesis follows from this section alone.

The arithmetic proof is human-readable and the finite part is replayed
with exact integer code.  This section is not a completed Lean
formalization.

\section{A corrected Selmer interface and a fixed-coefficient certificate}
\label{sec:sha-odd-bridge}

The earlier report \texttt{grh\_beyond\_dirichlet\_2026-11-25}, in
\texttt{sections/hecke\_selmer.tex}, identifies a finite-precision
quantifier obstruction. Its description matches the displayed theorem
in the compact edition, but it is not an obstruction to the readable
edition's endpoint. The distinction changes the proposed research
program: a full divisible Selmer bound is already present.

\subsection{What the readable edition actually proves}

In \texttt{New/Congruences/Selmer/Line.lean}, lines 120--144,
\texttt{divisibleSelmer} has coefficients
\(\operatorname{Frac}(\mathcal O')/\mathcal O'\), and
\texttt{TwistedLengthLe} bounds the length of this module for all
sufficiently large coefficient rings. The theorem
\texttt{exists\_anticyclotomic\_twist} in
\texttt{GeneralTheorem.lean}, lines 392--396, returns this property for
one character and all its first \(m\) powers. The intermediate precision
\(M\) is chosen with a strict margin above the error bound; it is not a
restriction of the final conclusion to coefficients modulo \(\ell^M\).

The key argument is explicit in
\texttt{Line.lean}, lines 84--95:
\path{length_selmerModule_divisibleQuot_le}. It is useful to expose
its elementary content.

\begin{proposition}[Saturation from one finite level]
\label{prop:one-level-saturation}
Let \(\mathcal O\) be a discrete valuation ring, let \(\pi\) be a
uniformizer, and let \(M\) be a torsion \(\mathcal O\)-module. If
\[
  \operatorname{length}_{\mathcal O}M[\pi^n]<n,
\]
then \(M=M[\pi^n]\), and in particular \(M\) has finite length.
If \(\mathcal O\) is the ring of integers of a finite extension of
\(\mathbb Q_\ell\), with absolute ramification index \(e_{\mathcal O}\),
put
\(\lambda_{\mathcal O}(N)=
\operatorname{length}_{\mathcal O}(N)/e_{\mathcal O}\).
Then
\[
 \lambda_{\mathcal O}\bigl(M[\ell^N]\bigr)\leq b<N
 \quad\Longrightarrow\quad
 M=M[\ell^N],\qquad \lambda_{\mathcal O}(M)\leq b.
\]
\end{proposition}

\begin{proof}
If an element of \(M\) has order \(\pi^r\) with \(r\geq n\), a suitable
multiple generates a submodule isomorphic to \(\mathcal O/\pi^n\)
inside \(M[\pi^n]\). This forces length at least \(n\).
Consequently every element is killed by \(\pi^{n-1}\), so the finite
submodule \(M[\pi^n]\) is all of \(M\).
For the normalized assertion use
\(\ell^N=u\pi^{e_{\mathcal O}N}\) for a unit \(u\).
\end{proof}

The strict inequality is sharp:
\(M=\operatorname{Frac}(\mathcal O)/\mathcal O\) satisfies
\(\lambda_{\mathcal O}(M[\ell^N])=N\) for every \(N\), but has infinite
length.

For the source's Selmer modules, the two algebraic identifications
needed here are proved in the repository. The module of locally
constant cocycles with divisible coefficients is torsion, because
a locally constant function on the compact absolute Galois group
has finite image and hence a common denominator
(\texttt{Line.lean}, lines 62--77). The module with coefficients
\(\mathcal O/a\) identifies with the \(a\)-torsion of the divisible
one (\texttt{Lib/GaloisRep/Main/DivisibleSelmerModules.lean},
lines 118--133). These are modules of cocycles satisfying fixed
linear relations, rather than arbitrary unrelated finite groups.

Thus even the compact theorem has a precise upgrade: retain all its
hypotheses and additionally suppose that its bound satisfies
\(B_i+\eta<M\) for every datum \(i\). After its single twist \(\psi\)
has been chosen, apply Proposition~\ref{prop:one-level-saturation}
to that same twist, coefficient ring, and local relations. The
conclusion is the corresponding full divisible Selmer bound for
\(\psi^d\), \(1\leq d\leq m\). No interchange of
\(\forall M\exists\psi\) and \(\exists\psi\forall M\) occurs.
The readable proof implements this step in
\texttt{twistedLengthLe\_pointChar}, lines 441--470, with
\(v_\ell(P(\zeta-1))<M\); its
\texttt{IsTorsionLevel}, lines 491--499, supplies the required margin.

This corrects one earlier obstruction. It does not identify the
source's local relations with the Bloch--Kato local conditions of an
arbitrary motive, remove its representation hypotheses, or turn a
small normalized length into zero.

\subsection{The exact odd-cocycle functor}

Here is a separate simplification of the remaining noncyclotomic
input. Fix a CM number field \(E\), a prime \(\ell\), a nonnegative
integer \(j\), and a chosen complex conjugation \(c\). Write
\(G_E=\operatorname{Gal}(\overline{\mathbb Q}/E)\), and let
\(\chi_\ell:G_E\longrightarrow\mathbb Z_\ell^\times\) be the cyclotomic
character. For a ring of integers \(\mathcal O\) of a finite extension
of \(\mathbb Q_\ell\), define
\[
 \mathcal A_{\ell,j,\nu}(E;\mathcal O)
\]
to be the \(\mathcal O\)-module of continuous functions
\(z:G_E\to\mathcal O/\ell^\nu\mathcal O\) satisfying:
\begin{enumerate}
\item \(z(gh)=z(g)+\chi_\ell(g)^{-j}z(h)\);
\item \(z\) vanishes on inertia at every finite place away from \(\ell\);
\item if \(j=0\), \(z\) vanishes on inertia above \(\ell\); if \(j>0\),
      \(z\) vanishes on the kernel of \(\chi_\ell\) in every
      decomposition group above \(\ell\);
\item if \(j>0\), \(z(cgc^{-1})=-z(g)\). For \(j=0\) this last
      condition is omitted.
\end{enumerate}
These are exactly the scalar-coordinate conditions of the source's
\texttt{shaOdd}, defined in
\path{Lib/GaloisRep/Main/SelmerModules.lean}, lines 174--190.
No quotient by coboundaries is taken. All local subsets in these
conditions are independent of the coefficient ring and of \(\nu\).

\begin{theorem}[Scalar extension for the source's strict odd cocycles]
\label{thm:sha-basechange}
For every finite coefficient extension \(\mathcal O/\mathbb Z_\ell\),
the natural map is an isomorphism:
\[
 \mathcal O\otimes_{\mathbb Z_\ell}
      \mathcal A_{\ell,j,\nu}(E;\mathbb Z_\ell)
 \ \xrightarrow{\ \sim\ }\
      \mathcal A_{\ell,j,\nu}(E;\mathcal O).
\]
If these modules have finite length, then
\[
 \operatorname{length}_{\mathcal O}
      \mathcal A_{\ell,j,\nu}(E;\mathcal O)
 =
 e_{\mathcal O}\,
 \operatorname{length}_{\mathbb Z_\ell}
      \mathcal A_{\ell,j,\nu}(E;\mathbb Z_\ell).
\]
Consequently their normalized lengths are independent of
\(\mathcal O\) and are nonnegative integers. The statements include
\(\ell=2\).
\end{theorem}

\begin{proof}
Choose a \(\mathbb Z_\ell\)-basis \(b_1,\ldots,b_d\) of \(\mathcal O\).
Reduction modulo the ideal \(\ell^\nu\mathcal O\), not modulo a fixed
power of a uniformizer, gives
\[
 \mathcal O/\ell^\nu\mathcal O
   =\bigoplus_{i=1}^d b_i(\mathbb Z_\ell/\ell^\nu).
\]
Write \(z=\sum_i b_i z_i\). The coordinate maps are maps between finite
discrete coefficient sets, so \(z\) is continuous if and only if all
\(z_i\) are continuous. Each defining equation has coefficients in
\(\mathbb Z_\ell\): this includes the cyclotomic action, both vanishing
conditions, and the equation defining oddness. Hence \(z\) satisfies
the conditions if and only if every \(z_i\) satisfies them. This proves
surjectivity and injectivity of the natural scalar-extension map.
No averaging operator \((1-c)/2\) is used, so no inversion of \(2\)
is required.

For finite length, write the base module as
\(\bigoplus_i\mathbb Z_\ell/\ell^{a_i}\).
After extension it is
\(\bigoplus_i\mathcal O/\ell^{a_i}\mathcal O\), whose length over
\(\mathcal O\) is \(e_{\mathcal O}\sum_i a_i\).
Finite length also descends: if the extended module has finite
length, its restriction to \(\mathbb Z_\ell\) is a finite-length module
containing the base module as a direct summand. The residue degree
does not enter the displayed composition-length factor.
\end{proof}

The earlier \texttt{ramification\_budget.tex} already proved the
generic invariance of normalized length under scalar extension and
identified fixed-module descent as a sufficient extra hypothesis.
Theorem~\ref{thm:sha-basechange} supplies the actual descent
identification for this particular source functor. It does not
supply a comparison from an untwisted Selmer group to a differently
twisted one.

\subsection{A single precision suffices for a proposed odd-cocycle bound}

Let \(\mathcal A_{\ell,j,\infty}(E)\) denote the same cocycle functor
with coefficients in the discrete module
\(\mathbb Q_\ell/\mathbb Z_\ell\).
A continuous function from \(G_E\) to this discrete module has finite
image, because \(G_E\) is compact. Its values therefore have a common
\(\ell\)-power denominator. The natural equivariant embedding
\(\mathbb Z_\ell/\ell^\nu\to\mathbb Q_\ell/\mathbb Z_\ell\),
\(\bar a\mapsto a/\ell^\nu\), gives
\[
 \mathcal A_{\ell,j,\nu}(E;\mathbb Z_\ell)
 \simeq \mathcal A_{\ell,j,\infty}(E)[\ell^\nu],
 \qquad
 \mathcal A_{\ell,j,\infty}(E)
   =\bigcup_{\nu\geq0}\mathcal A_{\ell,j,\infty}(E)[\ell^\nu].
\]
Every defining relation is preserved in both directions. In particular,
the strict local condition has not been replaced by a larger
cohomological condition.

\begin{samepage}
\begin{corollary}[One-precision certificate]
\label{cor:sha-one-precision}
For an integer \(B\geq0\), the following conditions are equivalent:
\begin{enumerate}
\item \(\operatorname{length}_{\mathbb Z_\ell}
        \mathcal A_{\ell,j,\nu}(E;\mathbb Z_\ell)\leq B\)
      for every \(\nu\);
\item \(\operatorname{length}_{\mathbb Z_\ell}
        \mathcal A_{\ell,j,B+1}(E;\mathbb Z_\ell)\leq B\);
\item \(\mathcal A_{\ell,j,\infty}(E)\) has length at most \(B\);
\item for every finite coefficient extension \(\mathcal O\) and every
      \(\nu\),
      \(\lambda_{\mathcal O}
          (\mathcal A_{\ell,j,\nu}(E;\mathcal O))\leq B\).
\end{enumerate}
In particular, the vanishing of the mod-\(\ell\) group is equivalent
to vanishing at every torsion level and over every finite coefficient
extension.
\end{corollary}
\end{samepage}

\begin{proof}
The implication from the first condition to the second is immediate.
The identification with torsion just proved and
Proposition~\ref{prop:one-level-saturation}, with \(N=B+1\), give the
third condition from the second. Taking torsion submodules gives the
first condition from the third.
Theorem~\ref{thm:sha-basechange} proves equivalence with the fourth.
\end{proof}

This is a fixed-field, fixed-exponent, fixed-local-condition
statement. It allows the coefficient ring to vary because a genuine
scalar-extension theorem has been proved. It allows the torsion level
to vary because all truncations belong to the same torsion module.
Changing the character, the representation, the field, or the local
relations requires a separate comparison.

\subsection{The exact remaining arithmetic target}

The source theorem \texttt{exists\_length\_shaOdd\_le}
(\texttt{Lib/Cyclotomic/OddSha.lean}, lines 645--653) requires
\(E\subseteq\mathbb Q(\mu_N)\). It constructs an integer
\[
  0<X\leq 2^{\,4(N+t+3)^{27}}h_E
\]
such that, for every prime \(\ell\), every allowed coefficient ring,
every \(j<t\), and every \(\nu\),
\(\operatorname{length}_{\mathcal O}\mathcal A_{\ell,j,\nu}
  \leq v_{\mathcal O}(X)\), with \(v_{\mathcal O}\) the
integer-valued discrete valuation.

For a noncyclotomic CM field, a proposed positive integer \(X\)
satisfies this entire bound if and only if
\[
 \boxed{\begin{gathered}
 \operatorname{length}_{\mathbb Z_\ell}
    \mathcal A_{\ell,j,v_\ell(X)+1}(E;\mathbb Z_\ell)
       \leq v_\ell(X)\\
 \text{for every prime }\ell\text{ and every }j<t.
 \end{gathered}}
\]
This follows from Corollary~\ref{cor:sha-one-precision}. The
coefficient-ring and all-precision quantifiers have been removed.
The infinitely many prime conditions have not: for every
\(\ell\nmid X\), the condition is mod-\(\ell\) vanishing.
Nor does this result give an algorithm that computes those Galois
cocycle groups.

Finally, the source uses the resulting odd-class cost through
\[
 \sum_\ell v_\ell(X)\log\ell=\log X.
\]
In \texttt{TraceRelations.lean}, lines 106--130, put
\(G=\texttt{chosenGroupOrder}\), \(h=\texttt{halfGap}\), and
\(f=\texttt{numInfPlaces}\). Its denominator error is
\[
 d_\ell=t\,v_\ell(X)
 +\bigl(8ft+2+(4t+1)h^{22}+h\bigr)v_\ell(G),
\]
and its clearing bound is
\[
 \beta_\ell=
   \operatorname{localDefect}_\ell
   +(1+2t)\operatorname{twistError}_\ell+d_\ell+v_\ell(G).
\]
The weighted sum, together with the \(U(1)\) class-number cost,
must fit within
\[
 \frac{\operatorname{errorUnit}}{mk}
   =\frac{fa(\log b)^3}{mk};
\]
see \texttt{FitTwistsExist.lean}, lines 232--260.
A replacement odd-cocycle theorem must therefore control \(\log X\)
uniformly along the varying fields and weights of a transfer.
Establishing a bound separately for each fixed field is insufficient
without this growth estimate.

The results in this section are proof extractions and explicit
interface clarifications. The saturation mechanism is already
formalized upstream; the scalar-extension theorem and its combined
certificate are proved here on paper. We make no claim that these
elementary module-theoretic facts are new in the literature, and we
have not completed a new Lean formalization of this combined interface.


\section{The remaining program beyond Dirichlet GRH}
\label{sec:program}

\subsection{A small nonabelian benchmark}

The splitting field of $X^3-X-1$ has Galois group $S_3$ and quadratic
subfield $K=\Q(\sqrt{-23})$. Its irreducible two-dimensional factor
is induced from a nontrivial cubic class-group character of $K$.
The familiar factorization of its degree-six Dedekind zeta function
contains the zeta factor, the quadratic Dirichlet factor, and the
square of this two-dimensional factor.  Thus zero control for the
remaining factor is a natural next target after Dirichlet GRH.
This is the benchmark developed in several incoming reports,
including \cite{prior-bridges,prior-synthesis}.

The combinatorial use of this field in the density theorem is
independent of its $L$-function. The coincidence of the fields allows
small arithmetic certificates to be shared; it does not make the
density theorem evidence for the new zero statement.

\paragraph{What the character result completes.}
One can choose tame characters with an allowed half of the split
support, a precursor trivial above $3$, and arbitrarily small
normalized inertia cost. A common residue field test certifies these
properties, and the finite conductor table quantifies their cost in
the specified class. This is stronger than merely constructing an
anticyclotomic character with the desired output properties.

\paragraph{What it does not complete.}
The character has not been shown to simultaneously meet every
representation-dependent exclusion, determinant nonvanishing
condition, small special-value ratio, and Selmer bound required by
the theta-lift construction. A character can satisfy all elementary
ramification restrictions and still be unusable for a given lift.
Nor does qualitative Chebotarev provide a conductor bound uniform in
all changing parameters at the scale needed by the proof. The table
shows exact finite costs, not an asymptotic conductor estimate.

\subsection{Four estimates that must refer to the same construction}

\begin{enumerate}
\item \textbf{A prime detector realized by divisibility.} The new
arithmetic coefficients must distinguish the required nonabelian
Frobenius data and arise in valuations of the actual derivative or
lifting construction. Replacing a Dirichlet weight by a class
function in a final sum supplies no such divisibility theorem.

\item \textbf{Short forms and enough independent images.} Both
bounds must hold for the replacement field, level, lattice, and
weights. The original proof uses almost-full spanning before it
absorbs the missing-rank error. A constant positive rank fraction by
itself need not preserve its error budget.

\item \textbf{A compatible lower determinant bound.} The lower
bound must concern the same integral structure and lift chain as the
upper bound. It must control all relevant local denominators and
the contribution of twists. Separate existence statements with
incompatible witnesses cannot be multiplied into a global bound.

\item \textbf{Uniform odd-cocycle and local losses.} The fixed-field
single-precision criterion in Section~\ref{sec:sha-odd-bridge}
removes redundant coefficient quantifiers. A transfer still needs
the actual bounds, including vanishing outside the prime support
of the proposed integer $X$, and a sufficiently small growth bound
for $\log X$ as the fields and weights vary.
\end{enumerate}

The final analytic zero criterion is relatively reusable once these
four inputs are proved. The prior corpus already develops many
versions using positive averages, multiplicative differences,
relative zeta functions, and group-theoretic detectors. Further
reformulation of that last step would add less than proving one of
the arithmetic estimates above.

\subsection{Algebraic K-theory: a concrete finite target with a global gap}

The reports on $K_{124}(\Z)$ and $K_{128}(\Z)$ suggest a second
credible route. This session replayed the supplied exact prefilter:
84 Bernoulli values agreed under two rational formulas and all 14
gcd certificates passed. The signed dimensions recorded for the
proposed spaces are $550,1240$ in degree $124$, and $605,1360$ in
degree $128$. These are reproducible finite preparations from
\cite{prior-ktheory}; they are not new vanishing theorems.

Three tasks remain essential: identify the source's strict cocycles
with the correct global cohomological object and local conditions;
make the representation and witness quantifiers uniform in the
relevant Hecke comparisons; and certify complete spectral coverage
and the required large determinants. A Bernoulli unit or a gcd equal
to one in a preliminary table cannot replace those tasks.

The source's full divisible Selmer theorem removes one overbroad
fixed-precision objection, but does not remove parity or local
condition mismatches. In particular a statement about a negative
Tate twist with strict local triviality is not automatically a
statement about the positive twists used in arithmetic $K$-theory.

\subsection{Why several larger targets are not immediate transfers}

\begin{center}\small
\begin{tabularx}{\textwidth}{@{}>{\raggedright\arraybackslash}p{0.23\textwidth}Y@{}}
\toprule
Target & Additional mathematics required beyond the established source endpoint\\
\midrule
BSD & A relation to the relevant elliptic-curve Selmer conditions,
ranks, and leading terms. The source's scalar-inertia and duality
assumptions are restrictive, and a bound for selected twists does
not imply the required assertion for the original curve.\\
Vandiver & The correct cyclotomic parity, treatment of the relevant
prime-primary classes, and uniformity in the varying cyclotomic
field. The source's strict odd module is not the conjecture itself.\\
Twin primes / binary Goldbach & Correlations of primes and the
associated sieve obstacles. Individual zero-free half-planes do not
supply these correlation estimates.\\
Subconvexity & Control for an individual automorphic form and its
analytic conductor. A determinant bound averaged over a large space
can hide the behavior of an individual member.\\
Fractional-ratio natural density above greedy & A mixed-sign
valuation construction. The explicit integral-ratio cube exchange
contains a forbidden mixed-step triple and cannot serve here.\\
\bottomrule
\end{tabularx}
\end{center}

These entries state what this report does not establish. They are not
claims that the techniques can never apply. For a further sustained
project, the most informative arithmetic milestone would be a
uniform strict odd-cocycle theorem for a specified noncyclotomic CM
tower together with a compatible prime-divisibility construction.
For a shorter combinatorial project, finite weighted optimization
for fractional-ratio ideal sets gives a precise certificate search
for the unresolved natural-density improvement.

\section{Verification, reproducibility, and limits}
\label{sec:validation}

The infinite mathematical assertions are established by the proofs in
this report. The accompanying programs check the finite constructions,
exact constants, and arithmetic certificates used by those proofs.
Their successful output is not a substitute for a proof of an
asymptotic density statement, a Galois-cohomology comparison, or a new
zero-location theorem.

\subsection{Frozen sources and the earlier-report corpus}

The source inspection uses \texttt{lsernau/grh-lean} at commit
\begin{center}\small
\texttt{72fead102f663d5c7a407964a13dc0cf9d977207}.
\end{center}
The archive contains 31 selected upstream source and documentation
files, including 25 Lean files. Their bytes were checked against the
Git blob SHA-1 identifiers returned for this commit, and SHA-256
digests and byte counts are recorded in
\path{sources/upstream_manifest.json}. The original license, notice,
and attribution files are retained. These checks establish the
identity of the inspected text.

The ProveIt corpus is frozen at commit
\texttt{5cc267af75a635199ac5d93a7d49cbce521daefc}.
All 46 relevant incoming ZIP archives, totalling 43,822,390 bytes,
were downloaded and checked against their Git blob identifiers;
their SHA-256 digests and sizes are recorded in
\path{data/archive_verification.json}. The inventory is in
\path{data/archive_inventory.csv}. Distinct filenames, including
numbered variants, were kept distinct. The review used the extracted
reports to identify previous results, correct the Selmer-precision
claim, and avoid presenting inherited finite computations as new.
Hash verification does not validate the mathematics inside an
earlier archive, and this review is not an independent reproof of
every claim in those reports.

The upstream audit is static. No complete build of either Lean
edition and no independent kernel replay were run in this
environment; Lean and Lake were not installed. The upstream
README reports its own build, axiom, and independent-checker results.
Those reports are distinguished from checks performed here.
The new combined statements in this report have human-readable
proofs; they are not described as newly Lean-checked.
The exact source anchors and audit limitations are retained in
\path{sources/audit_findings.json} and
\path{sources/upstream_audit.md}.

\subsection{Finite certificates and their scope}

\paragraph{Ideal-density constructions.}
The program \path{code/check_ideal_density.py} checks the
five-point exponent exchange, 10,000 ternary saturation instances,
five finite smooth-monoid examples, and the cube exchange in
dimensions \(2,3,4\). It also verifies explicit signed-step
counterexamples showing why the cube-exchange construction is not
valid for arbitrary fractional-ideal ratios.

For \(K=\mathbb Q(\sqrt{-23})\), the density-gap certificate uses
\(Y=10^7\). The prime-ideal product has 664,449 factors: two factors
for each of 332,112 split rational primes, 224 inert-prime-ideal
factors of norm \(p^2\), and one ramified factor. The product is
rounded downward after every multiplication using an integer
denominator \(10^{40}\), so the reported value is a rigorous lower
bound rather than a floating-point estimate. With the proved tail
bound \(2/Y\), the recorded rational gap exceeds \(1/20{,}000{,}000\).
The full output is \path{data/ideal_density_certificate.json}.

The separate program \path{code/verify_ideal_natural_gain.py}
independently verifies the rational arithmetic for the four selected
prime ideals of norms \(2,2,3,3\), the simple splitting of \(2\) and
\(3\) in the defining quadratic order, and the elementary zeta-value
upper bounds used in the proof. It certifies the strict natural-density
gain
\[
 \frac{16658080679}{13931406950400}>0.0011957.
\]
Its output is \path{data/ideal_natural_gain_certificate.json}.
This particular gain concerns proper integral-ideal ratios.
The all-number-fields extension and the natural-versus-upper-density
separation follow from the general arguments in the manuscript,
not from checking finitely many fields.

\paragraph{Auxiliary-character conductor minima.}
The program \path{code/check_source_twists.py} checks five minima
under the source-compatible precursor condition and six minima
under the weaker condition on the final anticyclotomic character.
It uses exact modular arithmetic and arithmetic-progression
Eratosthenes sieves. The progression sieve is cross-checked against
an ordinary sieve through \(10^4\); modular square roots are
cross-checked by exhaustive search at odd primes through \(500\).

The certificate includes the exclusion of every smaller permitted
multipair support using the reductions proved in
Section~\ref{sec:source_twists}. In the exceptional two-pair search,
five support pairs are tested with both orientations, and the
corresponding relation determinants are nonzero modulo \(3\).
The outputs are
\path{data/source_twist_certificate.json} and
\path{data/source_twist_minima.csv}.
The class-field interpretation, the scope of the admissible
characters, and the reduction to these finite searches are proved
in the text. The calculation does not establish the required
automorphic determinant estimates or a nonabelian RH.

\paragraph{Module-theoretic interfaces.}
The program \path{code/check_module_bridge.py} performs 46,332
one-level saturation checks and 50,193 finite-certificate equivalence
checks on explicitly specified torsion-module types. It also
enumerates 30 finite cyclic-cocycle scalar-extension models,
including oddness in residue characteristic \(2\).
The results are in \path{data/module_bridge_checks.json}.
These are exact finite sanity checks for
Section~\ref{sec:sha-odd-bridge}; they do not compute the arithmetic
groups \(\mathcal A_{\ell,j,\nu}(E;\mathbb Z_\ell)\).
The proof of scalar extension and the compactness argument are
supplied separately in that section.

\paragraph{Inherited \(K\)-theory preparation.}
The prior-report verifier \path{code/verify_ktheory_prefilter.py}
was replayed with exact rational arithmetic. For the 14 even
parameters \(54\leq e\leq80\), it compares an exponential-series
recurrence with the Bernoulli-polynomial formula in 84 evaluations,
checks the corresponding greatest common divisors and explicit
B\'ezout certificates, and computes eight compact invariant-space
dimensions. The record is
\path{data/ktheory_prefilter_rerun.json}.
In particular, the two displayed gcds for degrees \(124\) and
\(128\) are \(4\). These computations are inherited preparatory
certificates. They neither compute the missing Hecke determinants
nor certify a new algebraic \(K\)-group vanishing.

\subsection{Replaying the supplied computations}

The five programs use Python's standard library only. Run the
following commands, with assertions enabled, from the directory
containing \path{report.tex}; the supplied \path{data/} directory
receives the records. Python 3.10 or later is sufficient.

\begin{minipage}{\textwidth}
\begin{verbatim}
python3 code/check_ideal_density.py --cutoff 10000000 \
  --output data/ideal_density_certificate.json

python3 code/verify_ideal_natural_gain.py \
  --output data/ideal_natural_gain_certificate.json

python3 code/check_source_twists.py

python3 code/check_module_bridge.py \
  --output data/module_bridge_checks.json

python3 code/verify_ktheory_prefilter.py --ranks \
  --output data/ktheory_prefilter_rerun.json
\end{verbatim}
\end{minipage}

The character verifier writes both its JSON certificate and its CSV
table automatically. Verification logs are included under
\path{verification/}. Replaying these commands reproduces the
finite certificates; rebuilding the upstream Lean development
requires its separately pinned Lean and Mathlib environment.

The resulting evidence supports the distinctions used throughout
this report: proved general statements, exact finite certificates,
source theorem extractions, and conditional arithmetic transfers.
No finite prime search is used as evidence that all zeros lie on a
line, and no small normalized length is identified with vanishing
without the hypotheses that justify that inference.


\begin{thebibliography}{99}
\small\raggedright

\bibitem{repo}
The \texttt{grh-lean} contributors,
\emph{GRH for Dirichlet L-functions, in Lean 4},
research repository, 2026.
Revision \texttt{72fead102f663d5c7a407964a13dc0cf9d977207};
see the root README and \path{grh/src/Statements.lean}.
The README attributes the underlying proof to Levent Alp\"oge with
Claude and the shorter version to the repository owner, Anthropic
colleagues, and Claude.
\url{https://github.com/lsernau/grh-lean/tree/72fead102f663d5c7a407964a13dc0cf9d977207}.

\bibitem{guide}
The \texttt{grh-lean} contributors,
\emph{Reading guide}, readable edition of \emph{GRH for Dirichlet
L-functions, in Lean 4}, 2026, revision as in \cite{repo}.
\url{https://github.com/lsernau/grh-lean/blob/72fead102f663d5c7a407964a13dc0cf9d977207/grh/GUIDE.md}.

\bibitem{best}
Andrew Best, Karen Huan, Nathan McNew, Steven J. Miller,
Jasmine Powell, Kimsy Tor, and Madeleine Weinstein,
\emph{Geometric-progression-free sets over quadratic number fields},
Proceedings of the Royal Society of Edinburgh Section A: Mathematics
\textbf{147} (2017), no.~2, 245--262.
DOI: \href{https://doi.org/10.1017/S030821051600010X}
{10.1017/S030821051600010X}.
The open preprint contains the numbered questions used in this report:
\url{https://arxiv.org/abs/1412.0999}.

\bibitem{murty-vanorder}
M. Ram Murty and Jeanine Van Order,
\emph{Counting integral ideals in a number field},
Expositiones Mathematicae \textbf{25} (2007), no.~1, 53--66.
DOI: \href{https://doi.org/10.1016/j.exmath.2006.08.001}
{10.1016/j.exmath.2006.08.001}.
Author-hosted version:
\url{https://www.math.uni-bielefeld.de/vanorder/count.pdf}.

\bibitem{mcnew}
Nathan McNew,
\emph{On sets of integers which contain no three terms in geometric
progression},
Mathematics of Computation \textbf{84} (2015), 2893--2910.
DOI: \href{https://doi.org/10.1090/mcom/2979}{10.1090/mcom/2979}.
See especially Theorem~3.3 for an integer-ratio natural-density
improvement, Section~4 for finite-prime optimization, and Section~7
for the extremal-density questions.
\url{https://www.nathanmcnew.com/GPFsets.pdf}.

\bibitem{milneCFT}
J. S. Milne,
\emph{Class Field Theory}, course notes, version~4.03,
6 August 2020.
See Chapter~V for ray class groups and global reciprocity,
Chapter~VI for splitting densities, and Chapter~VIII, Section~7,
for the Chebotarev density theorem.
\url{https://www.jmilne.org/math/CourseNotes/CFT.pdf}.

\bibitem{sourceCMSplit}
The \texttt{grh-lean} contributors,
\path{grh/src/Lib/ClassField/Global/CMSplittingCharacters.lean},
readable edition, revision as in \cite{repo}.
See the documentation and the lemmas
\texttt{IsFromU1.eq\_one\_of\_mem\_decompGroup} and
\texttt{IsFromU1.eq\_one\_of\_mem\_inertiaAbove}.
\url{https://github.com/lsernau/grh-lean/blob/72fead102f663d5c7a407964a13dc0cf9d977207/grh/src/Lib/ClassField/Global/CMSplittingCharacters.lean}.

\bibitem{sourceAuxPrimes}
The \texttt{grh-lean} contributors,
\path{grh/src/New/Congruences/Selmer/AuxiliaryPrimes.lean},
readable edition, revision as in \cite{repo}.
See \texttt{IsAuxPrimeDataAt},
\texttt{exists\_anticyclotomicChar\_of\_frob\_fixes\_roots}, and
\texttt{exists\_isAuxPrimeDataAt}.
\url{https://github.com/lsernau/grh-lean/blob/72fead102f663d5c7a407964a13dc0cf9d977207/grh/src/New/Congruences/Selmer/AuxiliaryPrimes.lean}.

\bibitem{incoming}
Vladimir Reshetnikov, repository maintainer,
\emph{ProveIt: incoming research reports},
revision \texttt{5cc267af75a635199ac5d93a7d49cbce521daefc},
inspected 26 November 2026.
Archive-specific references below identify research artifacts;
their inclusion does not imply independent verification of every claim.
\url{https://github.com/VladimirReshetnikov/ProveIt/tree/5cc267af75a635199ac5d93a7d49cbce521daefc/docs/incoming}.

\bibitem{prior-bridges}
\emph{Beyond Dirichlet GRH: Arithmetic twists, sparse equations,
and the remaining transfer barriers},
research report prepared for Vladimir Reshetnikov,
25 November 2026.
Archive \path{grh_transfer_bridges_2026-11-25.zip} in \cite{incoming};
especially Sections~3--5.
\url{https://github.com/VladimirReshetnikov/ProveIt/blob/5cc267af75a635199ac5d93a7d49cbce521daefc/docs/incoming/grh_transfer_bridges_2026-11-25.zip}.

\bibitem{prior-synthesis}
\emph{Beyond Dirichlet GRH: Validated arithmetic and transferable
proof machinery},
research report prepared for Vladimir Reshetnikov,
25 November 2026.
Archive \path{grh_transfer_synthesis_2026-11-25.zip} in
\cite{incoming}.
\url{https://github.com/VladimirReshetnikov/ProveIt/blob/5cc267af75a635199ac5d93a7d49cbce521daefc/docs/incoming/grh_transfer_synthesis_2026-11-25.zip}.

\bibitem{prior-ktheory}
\emph{Beyond Dirichlet GRH: Averaged zero detection and exact
K-theory transfer tests},
research report prepared for Vladimir Reshetnikov,
25 November 2026.
Archive \path{grh_method_transfers_2026-11-25 (2).zip} in
\cite{incoming}; the present report replays its exact Bernoulli and
gcd prefilter, with the recorded signed dimensions.
\url{https://github.com/VladimirReshetnikov/ProveIt/blob/5cc267af75a635199ac5d93a7d49cbce521daefc/docs/incoming/grh_method_transfers_2026-11-25%20(2).zip}.

\bibitem{prior-density-gaps}
\emph{Geometric progressions and nonabelian zero tests:
Three transfers from the Lean proof of Dirichlet GRH},
research note, 25 November 2026, incoming research package
\path{grh_transfers_density_gaps_and_zero_tests.zip} in
\cite{incoming}.
The relevant arguments are in \path{integers.tex},
\path{integer_more.tex}, \path{natural_density.tex}, and
\path{gap.tex} inside the package.
\url{https://github.com/VladimirReshetnikov/ProveIt/blob/5cc267af75a635199ac5d93a7d49cbce521daefc/docs/incoming/grh_transfers_density_gaps_and_zero_tests.zip}.

\end{thebibliography}

\end{document}
